\section{Task Definition and Evaluation Protocol}
\label{sec:task}

\subsection{Unit of prediction and metric}
The unit of prediction is a (figure, sentence) pair. For each figure $F$ a system outputs a set of sentences
$S_F$, which is compared with the gold set $G_F$:
\[
P_F=\frac{|S_F\cap G_F|}{|S_F|},\qquad R_F=\frac{|S_F\cap G_F|}{|G_F|},\qquad
F1_F=\frac{2P_FR_F}{P_F+R_F}.
\]
Precision, recall and F1 are then averaged over figures, so that every figure counts equally. This is the
protocol of \citet{ramesh2015figsum}, which makes our numbers directly comparable with theirs. The selected set
includes the explicit sentences. ``F1'' without qualification always means this per-figure mean.

\subsection{Vocabulary}
\begin{table}[h]
\centering
\small
\begin{tabularx}{\linewidth}{>{\raggedright\arraybackslash\bfseries}p{3.6cm}L}
\toprule
\textnormal{\textbf{Term}} & \textbf{Meaning}\\
\midrule
Explicit (citing) sentence & A sentence whose text names figure $F$, as detected by our grammar or by the
citation tags of the source XML.\\
Near implicit sentence & A sentence that does not name $F$ and lies within three sentences of a sentence that
does.\\
Far implicit sentence & A sentence that does not name $F$ and lies more than three sentences from every
citation of $F$ (or $F$ is never cited).\\
Out-of-fold (OOF) score & A score given to a paper by a model that was not trained on that paper. The training
papers are split into five folds of three papers; a model trained on four folds scores the fifth, and rotating
gives every paper an out-of-fold score.\\
AUC & The probability that a randomly chosen gold pair receives a higher score than a randomly chosen non-gold
pair: 0.5 for random scores, 1.0 for a perfect ranking. It measures ranking quality independently of any
cut-off. Example: gold scores $\{0.9, 0.4\}$ and non-gold scores $\{0.5, 0.2, 0.1\}$ give five wins in six
comparisons, so AUC $=5/6$. \emph{Near} and \emph{far} AUC restrict the comparison to near and far pairs.\\
\bottomrule
\end{tabularx}
\caption{Terms used throughout the report.}
\label{tab:vocab}
\end{table}

When we break results down by category, the explicit, near and far columns are computed from pooled sentence
counts (correct / selected / gold over all figures), while the overall P, R and F1 are per-figure means.

\subsection{Protocol and safeguards}
\label{sec:protocol}
The benchmark is small, so we took particular care that reported numbers are not flattered by tuning.
\begin{itemize}
  \item \textbf{A fixed split.} Fifteen papers form the training pool and three papers form a test set
        (Section~\ref{sec:split}). The split was fixed before the cross-encoder and
        span-tagger runs of Section~\ref{sec:models} were trained. Earlier exploratory phases used a 13-paper
        protocol; results from those phases are marked as such.
  \item \textbf{Everything out-of-fold.} All model selection uses out-of-fold scores on the training pool.
        Each configuration is trained with three random seeds per fold and the three scores of a pair are
        averaged. Even the decision threshold is chosen out-of-fold: the threshold applied to a fold is the one
        that maximises F1 on the other four folds.
  \item \textbf{A sealed, pre-registered test.} Before any test score was computed we wrote down the primary
        model, its decision rule, the source of its threshold and the criterion for success
        (file \texttt{TEST\_PLAN.md} in the repository). The test set was then scored once.
\end{itemize}
